% Français 9115, batch 37 — Gallica views 721–740.
% Pass: Opus 5.5 (claude-opus-5-5), 2026-09-22 — first pass, unchecked against the views by a human.
%
% What the twenty views hold. Two different things, one after the other.
% Views 721–733: number theory, quadratic residues, in her fast working
% hand, on small leaves and slips mounted two to a guard. Views 734–740 (and
% on past the batch): fair extracts, in a regular hand, from memoirs of the
% Turin Academy on sound and vibrating strings (Lagrange, Euler), with her
% own remarks in square brackets. Every view was read from the local
% mirror; none is missing. Fourteen views are transcribed, written backs
% included: 722, 723, 724, 726, 728, 729, 730, 734, 735, 736, 737, 738,
% 739, 740.
%
%   v. 722    (two slips, ink 357 above, 358 below; written on both
%             sides) residues. Slip 357: every number 3k−1 is a
%             non-residue of some prime, via p + 3aa and −3 non-residue
%             of the primes 3k−1 — « mais je ne sais pas le demontrer ».
%             Slip 358: « P. 125 L'auteur se borne a démontrer le cas des
%             nombres premiers dela forme 4n+1 … (N.o 139, § 146) »; the
%             general proposition « Tout nombre, excepté les quarrés
%             positivement pris, est non residu de quelque nombre
%             premier », reduced to 4n+2 and 4n+3; the case 4n+2 done,
%             « Je n'ai pas vu comment demontrer −(4n+2) non residu »;
%             « Retournez ».
%   v. 723    their backs: three lines on 357 (numbers both 3k−1 and
%             8n+1, i.e. 24n+17); on 358, the cases +(4n+3) and −(4n+3),
%             ending « Tout nombre dela forme 4k+3 pris tant avec le signe
%             + qu'avec le signe − » at the foot of the slip.
%   v. 724    (ink 359, headed « §126 ») a prime 8n+1 is a non-residue of
%             a smaller prime: the case 24k+17 (as 12k+5, via p − 3aa and
%             « 3 non residu mod. 12k+5 et 12k+7, (V. P. 116) ») complete;
%             the case 24k+1 begun as a sum of two squares and residues
%             mod 5, stopping on a congruence. Fairer hand.
%   v. 726    (ink 361 above, 360 below) 5 and −3 as non-residues of the
%             factors of numbers of a form F; an induction « Supposons que
%             tous les nombres premiers, jusqu'a P, pour lesquels a est
%             residu, soient delaforme ma + b² ».
%   v. 728    (ink 363 above, cancelled by heavy vertical strokes; 362
%             below) « P. 116. »: for primes 4k+3, n non-residue ⇒ −n
%             residue, by a primitive root; « L'induction a mené l'auteur a
%             établir pour les nombres 12n+11 3 residu et −3 non residu »;
%             and six lines of algebra on A = ma + b², a = d² − nA.
%   v. 729    (photographed sideways) the backs of 363 and, probably, 362:
%             « Pour tout nombre premier 12n+5 dans lesquels 3n+1 est
%             premier 3 est racine primitive », with examples 29, 53 and
%             two struck ones (77, 93); and a cut pencil fragment in
%             another, larger hand on determinants (« la déterminante … de
%             l'ordre n+1 »), lines cut at both edges.
%   v. 730    (ink 364) « Nombres premiers delaforme 2^{2^i}+1 »: the
%             smaller Fermat primes are non-residues, hence primitive
%             roots, of a larger one — an argument she then cancels
%             between the lines, « cela est faux parcequ'il suffit que
%             … ≡ x² »; and 3 and 7 non-residues mod 2^{2^i}+1, « c'est
%             cequej'avois deja trouvé d'une autre maniere ». Her « R »
%             between two letters (a R a′) means « is a residue of »; the
%             « T. » before it is not explained on the leaf.
%   v. 734    (ink 365, a large leaf) « Memoire de turin [2d ?] année 1761
%             Pag. 16 Theorie du son La Grange »: against Newton's curve EKS
%             (Principia II, prop. 49), which cannot be a circle, nor any
%             algebraic or transcendental curve; φ(t − xT/√(2hA)).
%   v. 735    its back: the need for « fonctions irrégulières et
%             discontinues »; « Probleme »: an infinity of mobile points,
%             d²z/dt² = c d²z/dx², multiplied by M dx and integrated by
%             parts.
%   v. 736–737 (ink 366, both sides) the solution continued: d²M/dx² = KM,
%             s = ∫zM dx, s and r in exponentials, then in cos and sin of
%             t√(−cK); M = A sin x√(−K), a√(−K) = νπ/a; « Remarque »: each
%             point must be treated as isolated, an infinity of values of
%             K.
%   v. 738–739 (ink 367, both sides) « La Grange (1 2d. de T. P. 69) »:
%             the sum of cos x + cos 2x + … = −½ via imaginary
%             exponentials; d'Alembert's objection (the last term
%             e^{∞√−1} is not zero) and Lagrange's reply (« S. 2 T. P.
%             317 »), with her own remark in brackets — x^∞ is only the
%             limit of x, x², x³ …; d'Alembert's example x = 45°, quoted
%             with repeated guillemets, and her bracketed answer that
%             cos mx − cos(m+1)x over 2(1 − cos x) takes four values.
%   v. 740    (ink 368) « 3 V. Des Memoires de Turin — Eclercissement sur
%             le mouvement des cordes vibrantes Par Euler », « Question »,
%             a pen figure (axis AB, curves AMB and ASB, ordinate MXV);
%             runs on to its back, v. 741, past the batch.
%
% Views not transcribed. v. 721: blank backs of the two slips inked 355
% and 356 (v. 720, previous batch), their text showing through reversed; a
% round library stamp shows faintly on the upper one. v. 725: blank back of
% 359. v. 727: blank backs of 361 and 360, 360's text showing through
% strongly. v. 731: blank back of 364, the fold marks of v. 730 visible.
% v. 732: a leaf blank on this side, with no ink number; v. 733: its other
% side, blank (both checked at raised contrast). No repeat, no béquet, no
% microfilm target.
%
% Seams. At 720/721: v. 720 (previous batch, glanced at only) carries two
% slips inked 355 and 356 on the same subject, residues — 356 reads « je me
% propose de demontrer qu'un nombre quelconque dela forme 4k+3 … tant avec
% le signe + qu'avec le signe − est non residu de quelque nombre premier
% moindre que lui », which is the proposition v. 723 announces in its last
% line. v. 721 is blank, so no sentence is cut at 720/721. At 740/741: the
% Euler extract of v. 740 stops mid-sentence (« … tout comme les deux
% conditions prescrites l'exigent : ») and continues on the back, v. 741.
%
% The numbers on the leaves, and why no \folio{} is written. (1) The ink
% series at the top right, in the larger rounder hand: 357 and 358 (v. 722),
% 359 (724), 361 above 360 (726), 363 above 362 (728), 364 (730), 365 (734),
% 366 (736), 367 (738), 368 (740, its top also showing at the head of
% v. 738 above that leaf). Never on the backs. The leaf at v. 732–733
% carries none, so the series skips it. On two guards the upper slip bears
% the higher number (361/360, 363/362). Penned, as everywhere in this
% volume: no \folio{}. Its 5 has a long descending tail and reads like a 9
% at band scale; the 9 of 359 is looped and closed. (2) Her own references,
% in the text: « P. 125 », « N.o 139 », « § 146 » (v. 722), « §126 » at
% the head of v. 724, « (V. P. 116) » (v. 724), « P. 116. » at the head of
% v. 728. As the pass's own inference, not a reading: 125, 126, 139, 146
% and the propositions quoted next to them match the article numbering of
% Gauss's Disquisitiones arithmeticae (« L'auteur … son théorème
% fondamental »), and « P. » may be a page of a French translation or an
% abbreviation of article; nothing on the leaves names the book. (3)
% Printed-source references in the Turin extracts: « année 1761 Pag. 16 »
% (v. 734), « Prop. 49 liv. 2 des Principes » (v. 734), « 1 2d. de T.
% P. 69 » and « S. 2 T. P. 317 » (v. 738), « 3 V. » (v. 740). Every
% number above was read on its view; none is inferred.
%
% Manuscript C. None of these leaves is Laubenbacher and Pengelley's
% « Manuscript C » (ff. 348r–349r): no Fermat equation, no even exponents.
% The number theory here is quadratic residues and Fermat primes, inked
% 357–364. At one number per leaf, blank leaves included, 348–349 would
% fall about eighteen views before v. 722, around v. 704 (batch 36) — an
% estimate from the series, not a reading; v. 720 already carries 355–356,
% also on residues.
%
% Notation. Hers, kept. The congruence sign is written with three strokes
% and is set ≡ wherever it has three (distinct from her =). « sin »,
% « cos », « log. » are set \text{}. Her « R » between two numbers (a R a′,
% « a residu de a′ ») is kept as a capital. In the Turin extracts: $s$ for a
% long-stemmed letter; $\mu$ and $\nu$ for a hooked u and a looped v; $\alpha$
% for a small u in XS = αXM (v. 740). The page's own square brackets are
% kept. Words she underlines are set \emph{}. A long s is transcribed s.
% Spelling is the page's and is not flagged: « bornoit », « dela »,
% « parconséquent », « quarrés », « fesons », « s'agitte », « tems »,
% « evanoüisse », « appeller », « Eclercissement »,
% « j'avois ». Where the leaf and the calculation disagree the leaf stands,
% with a note (v. 724: 25n for 25n²; v. 736: the exponent t−√cK; v. 737:
% the two equations after substitution, and « sin. x√−1 »).
%
% What could not be read, and what is offered with doubt. \ill{} stands
% in the text: the second member of a congruence under a cancelling
% stroke (v. 728); a struck calculation and, in the pencil fragment, the
% letters cut by its edges, thirteen in all (v. 729); a short struck word
% before « sauroit » (v. 734); the base and sign of an exponential under a
% blot (v. 736); a struck group after « a moins que » (v. 739). A small
% struck group between two lines of v. 722 is described in its note.
% \uncertain{} stands for « P. », « V. », « § », « est » (v. 722); the 2
% of 24n+17 (v. 723); the denominator 1 of a − 1 (v. 728); « est »,
% « Jaune », « Ce Mon… » (v. 729); « 2d ? » and the u of t − ux (v. 734);
% « ont » and the denominator a (v. 737); « 1 », « P. », « S. », « y »
% (v. 738); « Pag. D. » (v. 739); « V. » (v. 740). Nothing was guessed to
% fill a gap.
%
% Nothing here was checked against any published edition of Sophie
% Germain's papers, of Gauss, or of the Turin memoirs.
\documentclass[11pt,a4paper]{article}
% The preamble shared by every transcription of the mathematicians' manuscripts in Gallica.
%
% It rests on one idea: **uncertainty compounds, it does not resolve.** A
% transcription that smooths over a doubtful word has destroyed the only thing
% it was made for — a reader must be able to tell, at every point, what was on
% the page and what was supplied. The apparatus macros below are therefore the
% critical apparatus, not decoration.
%
% The same commands are recognised by `scripts/render.mjs`, which produces the
% site's reading view, and by `scripts/tei.mjs`. A macro added here must be
% added there too, or rendering fails loudly — which is the wanted behaviour.
%
% Comments here are English, like the rest of the repository. The documents
% this preamble sets are French, and that is deliberate: see the skills.

\ProvidesPackage{germain}[2026/09/15 Transcriptions of mathematicians' manuscripts in Gallica]

\usepackage{amsmath,amssymb,amsthm}
\usepackage{mathrsfs}
% Kept from the parent preamble so the renderer's diagram support stays
% available; these manuscripts draw no commutative diagrams, and nothing here uses it.
\usepackage{tikz-cd}
\usepackage[english,french]{babel}
\usepackage{fontspec}

% --- Greek ------------------------------------------------------------------
% The text face stays fontspec's default, Latin Modern, which has no Greek:
% XeTeX drops every glyph it lacks with a line in the log and nothing on the
% page, and the Greek words the manuscripts quote vanished from the PDFs. So
% the Greek and Greek Extended blocks get a XeTeX character class of their own,
% and entering or leaving a run of it switches the family to CMU Serif — the
% same Computer Modern design, with polytonic Greek — and back. Only the
% family changes: a Greek word in \textit or \textbf keeps its shape and series.
% The transitions to and from every other class carry the tokens that class
% has towards an ordinary letter, so French spacing before « ; » or inside
% « » still holds after a Greek word. No grouping, so no brace or \endgroup
% can come out unbalanced; maths is left alone, since XeTeX inserts nothing
% there. Kept identical in `exercices.sty`.
\makeatletter
\newfontfamily\germainGreekfont{cmun}[
  Extension=.otf, NFSSFamily=germaingreek,
  UprightFont=*rm, ItalicFont=*ti, SlantedFont=*sl,
  BoldFont=*bx, BoldItalicFont=*bi, BoldSlantedFont=*bl]
\newXeTeXintercharclass\germain@greek
\def\germain@greekrange#1#2{%
  \@tempcnta=#1\relax
  \loop
    \XeTeXcharclass\@tempcnta=\germain@greek
    \ifnum\@tempcnta<#2\advance\@tempcnta\@ne\repeat}
\germain@greekrange{"0370}{"03FF}
\germain@greekrange{"1F00}{"1FFF}
\def\germain@greekprev{\rmdefault}
\def\germain@greekin{%
  \edef\germain@greekprev{\f@family}\fontfamily{germaingreek}\selectfont}
\def\germain@greekout{\fontfamily{\germain@greekprev}\selectfont}
\def\germain@greekjoin{%
  \XeTeXinterchartokenstate=\@ne
  \@tempcnta=\z@
  \loop
    \ifnum\@tempcnta=\germain@greek\else
      \XeTeXinterchartoks\germain@greek\@tempcnta=\expandafter{%
        \expandafter\germain@greekout\the\XeTeXinterchartoks\z@\@tempcnta}%
      \XeTeXinterchartoks\@tempcnta\germain@greek=\expandafter{%
        \the\XeTeXinterchartoks\@tempcnta\z@\germain@greekin}%
    \fi
    \ifnum\@tempcnta<4095\advance\@tempcnta\@ne\repeat}
% babel-french sets its own classes, and saves and restores the interchar
% state, each time a language is selected: join again after it does.
\addto\extrasfrench{\germain@greekjoin}
\addto\extrasenglish{\germain@greekjoin}
\AtBeginDocument{\germain@greekjoin}
\makeatother

\usepackage{geometry}
\geometry{a4paper,margin=25mm,marginparwidth=28mm}
\usepackage{marginnote}
\usepackage{xcolor}
\usepackage[normalem]{ulem}
\setlength{\emergencystretch}{3em}
\usepackage{hyperref}
\hypersetup{colorlinks=true,linkcolor=black,urlcolor=black,pdfborder={0 0 0}}

% --- Batch metadata ---------------------------------------------------------
% Taken as they stand by the rendering script, which shows them at the head of
% the reading view. They fix what the file claims to cover. The macro names are
% the parent project's, so that the scripts stay shared; \folder here means the
% volume's slug — `fr-9115` — and \pages counts Gallica views.

\newcommand{\folder}[1]{\gdef\grothFolder{#1}}
\newcommand{\batch}[1]{\gdef\grothBatch{#1}}
\newcommand{\pages}[2]{\gdef\grothFirst{#1}\gdef\grothLast{#2}}
\newcommand{\foldertitle}[1]{\gdef\grothFoldertitle{#1}}

% The holder's own shelfmark, printed as they write it: « Français 9115 ».
\newcommand{\shelfmark}[1]{\gdef\grothShelfmark{#1}}

% The dating, copied verbatim from the holder's catalogue — never inferred
% here. « XIXe siècle » is the BnF's; a pass that dates a leaf from its content
% says so in a \note{}, as its own inference.
\newcommand{\dating}[1]{\gdef\grothDating{#1}}

% The status notice, printed in the title block and repeated as a diagonal
% watermark on every page of the PDF. The manuscripts are in the public
% domain; what this line declares is not a right but a quality — an
% unverified first pass by a machine. The skills write the line; a file
% without it gets no watermark, so leaving it out is visible in the output.
\newcommand{\watermark}[1]{\gdef\grothWatermark{#1}}

\gdef\grothFolder{?}\gdef\grothBatch{}\gdef\grothFirst{?}\gdef\grothLast{?}%
\gdef\grothFoldertitle{}\gdef\grothShelfmark{}\gdef\grothDating{}\gdef\grothWatermark{}

% --- The résumé -------------------------------------------------------------
\newenvironment{resume}
  {\small\setlength{\parskip}{0.45em}}
  {\par\medskip\hrule\medskip}

% --- Keywords ---------------------------------------------------------------
% In English, closing the modernised reading's résumé: the modern vocabulary
% under which to search for what the leaves construct. The manifest extracts
% them to tag the volume — this line is the single source of the tags.
\newcommand{\keywords}[1]{\par\smallskip\noindent
  {\footnotesize\textsc{Keywords} --- \textit{#1}}\par}

% --- The critical apparatus -------------------------------------------------

% The Gallica view number, in the margin. This is the anchor that lets a
% disagreement over a formula be settled by looking at *one* image rather than
% a volume of 750. The site uses it to turn the facsimile as the transcription
% is scrolled. A view is one scanned image: usually one side of a leaf.
\newcommand{\page}[1]{%
  \marginnote{\footnotesize\color{gray!70}\textsf{v.\,#1}}%
  \label{page:#1}\ignorespaces}

% The BnF's pencilled foliation, where the leaf carries it and a pass can read
% it: « 198r ». This is what the literature cites — « ff. 198r–208v » — and
% the one line that turns a citation into a view. Recorded, never inferred:
% a folio a pass cannot read is not written.
\newcommand{\folio}[1]{%
  \marginnote{\footnotesize\color{gray!50}\textsf{f.\,#1}}\ignorespaces}

% The modernised reading's coarser counterpart to \page{}: which views a
% section is drawn from.
\newcommand{\pagerange}[2]{%
  \marginnote{\footnotesize\color{gray!70}\textsf{v.\,#1--#2}}%
  \ignorespaces}

% Illegible: never guessed.
\newcommand{\ill}{\textcolor{red!60!black}{[\,\ldots\,]}}

% A doubtful reading: the word is given, and flagged as uncertain.
\newcommand{\uncertain}[1]{\uline{#1}}

% An editorial addition — a missing letter, an implied word.
\newcommand{\add}[1]{\textcolor{blue!50!black}{[#1]}}

% Struck out by the writer. What was crossed out is part of the document.
\newcommand{\struck}[1]{\sout{#1}}

% The transcriber's note, on the state of the page or the mathematics.
\newcommand{\note}[1]{\par\smallskip{\small\color{gray!60!black}\textit{[#1]}}\par\smallskip}

% A marginal note of hers, distinct from the body of the text.
\newcommand{\marginal}[1]{\par\smallskip\noindent\hspace{1em}%
  {\small\color{gray!60!black}#1}\par\smallskip}

% --- Title ------------------------------------------------------------------
% The title block is French, like the documents themselves.

\AddToHook{shipout/background}{%
  \ifx\grothWatermark\empty\else
    \begin{tikzpicture}[remember picture, overlay]
      \node[rotate=52, scale=3.4, align=center, text=gray!18,
            font=\sffamily\bfseries]
        at (current page.center) {\grothWatermark};
    \end{tikzpicture}%
  \fi}

\AtBeginDocument{%
  \begin{center}
    {\large\bfseries Manuscrits de mathématiciens (Gallica) ---
      \ifx\grothShelfmark\empty\grothFolder\else\grothShelfmark\fi}\\[2pt]
    {\itshape\grothFoldertitle}\\[4pt]
    {\small vues \grothFirst--\grothLast
      \ifx\grothBatch\empty\else\ (lot \grothBatch)\fi}\\[2pt]
    \ifx\grothDating\empty\else
      {\small\color{gray!60!black}Datation du catalogue : \grothDating}\\[2pt]
    \fi
    \ifx\grothWatermark\empty\else
      {\small\bfseries\color{red!45!gray}\grothWatermark}\\[2pt]
    \fi
    {\footnotesize\color{gray!60!black}Transcription établie d'après les images IIIF de Gallica.
     Source gallica.bnf.fr / Bibliothèque nationale de France.\\
     Les lectures incertaines sont soulignées, les illisibles notées [\,\ldots\,],
     les ajouts entre crochets ; \textsf{v.} numérote les vues de Gallica,
     \textsf{f.} la foliotation au crayon de la BnF.}
  \end{center}
  \vspace{1em}}

\endinput


\folder{fr-9115}
\shelfmark{Français 9115}
\batch{37}
\pages{721}{740}
\dating{1801-1900}
\watermark{Lecture automatique\\non vérifiée}
\foldertitle{Papiers de Sophie GERMAIN. — Recueil de dissertations et problèmes mathématiques et physiques.}

\begin{document}

\note{Numérotation des feuillets. En haut à droite, à l'encre, d'une écriture plus ronde, la série qui court dans tout le volume : 357 et 358 (vue 722), 359 (vue 724), 361 au-dessus de 360 (vue 726), 363 au-dessus de 362 (vue 728), 364 (vue 730), 365 (vue 734), 366 (vue 736), 367 (vue 738), 368 (vue 740) ; jamais au dos. Le feuillet des vues 732--733, blanc, n'en porte pas. Cette série occupe la place de la foliotation de la Bibliothèque, mais elle est tracée à la plume et non au crayon, et aucun « f. » n'est donc porté. Tous ces nombres ont été lus sur la vue ; aucun n'est déduit. Ne sont pas transcrites : la vue 721, dos blancs des feuillets 355 et 356 ; la vue 725, dos blanc du feuillet 359 ; la vue 727, dos blancs des feuillets 361 et 360 ; la vue 731, dos blanc du feuillet 364 ; les vues 732 et 733, les deux côtés d'un feuillet blanc sans nombre.}

\page{722}

\note{Deux petits feuillets montés l'un au-dessus de l'autre sur la même feuille de garde, écrits des deux côtés ; leurs dos forment la vue 723. En haut à droite de chacun, à l'encre, le nombre de la série du volume : « 357 » sur le feuillet du haut, « 358 » sur celui du bas. Le 5 a la queue descendante qui le fait prendre pour un 9 ; on le distingue du 9 bouclé de « 359 » (vue 724).}

\marginal{\emph{retournez}}
Si on se bornoit a vouloir demontrer que tout nombre de la forme
$3k-1$ est non residu d'un nombre premier, sans exiger que ce nombre
soit $<$ que $3k-1$ ; on y parviendroit en considerant un nombre
de la forme $p+3aa$ dans lequel $p=3k-1$ ; car, $p+3aa$ est dela
forme $3k-1$ et parconséquent il est, ou premier, ou divisible par
un nombre de la même forme ; soit $q$ ce nombre on a
$p\equiv-3aa$ (mod.~$q$) mais $-3$ est non residu pour tout nombre de
la forme $3k-1$ donc $p$ est aussi non residu (mod.~$q$)
si $p$ est pair \emph{a} doit être supposé impair, et pair si $p$ est impair
afin que $p+3aa$ soit un nombre impair. car si ce nombre étoit
pair il seroit possible que $2$ fut le seul nombre $3k-1$ qui le divisât
\emph{et on} sait que \add{tout nombre impair $\equiv1$ (mod~2)} \struck{tout nombre est divisible} \struck{$-3$ est residu mod 2. car tous}
On voit que dans la suite des differentes valeurs que l'on peut donner a \emph{a}
il doit se trouver des valeurs de $p+3aa$ divisibles par des valeurs de $q<p$.
\marginal{mais je ne sais pas le demontrer}
\note{En tête du feuillet, à gauche, au-dessus de la première ligne, un mot souligné, lu « retournez » : la première lettre est peu formée, mais le même mot, souligné et bien lisible, est écrit au bas du feuillet inférieur ; le dos de ce feuillet-ci (vue 723) porte trois lignes. Le « est » de la deuxième ligne est repassé et empâté, non barré. Entre la troisième et la quatrième ligne, au-dessus de « un nombre », un petit groupe biffé et empâté qui commence par « $1+$ » : \ill{}. Le signe de « $p\equiv-3aa$ » a trois traits, nettement distincts du signe $=$ de « $p=3k-1$ » ; il est rendu $\equiv$, ici et partout où il a trois traits. « \emph{et on} » est souligné d'un trait épais, qui peut aussi être une rature ; la ligne suivante commence par « On voit que » et le sens ne dit pas si « et on sait que » est maintenu. « tout nombre impair $\equiv1$ (mod 2) » est écrit au-dessus de « tout nombre est divisible », qu'il remplace ; la fin de la ligne, « $-3$ est residu mod 2. car tous », est barrée d'un trait épais qui court jusqu'au bord. « mais je ne sais pas le demontrer » est écrit en trois lignes serrées au bout de la dernière ligne, contre le bord droit du feuillet.}

\uncertain{P.}~125 L'auteur se borne a démontrer le \struck{tout} cas des nombres premiers
dela forme $4n+1$ parcequ'il n'emploie que celui là dans
sa demonstration de son théorème \struck{fondamental} (\uncertain{V.} N.o 139. \uncertain{§} 146)
si on vouloit demontrer la proposition generale :

Tout nombre, excepté les quarrés positivement pris, \uncertain{est} \struck{résidu}
non residu de quelque nombre premier.

Nous observerons qu'en supposant la chose accordée pour les nombres
$4n+1$ il ne reste a examiner que ceux $4n+2$, $4n+3$ parceque
les nombres $4n$ étant multiples d'un quarré seront non residu
si $n$ l'est.

Pour les nombres, $4n+2$

Soit $4n+2+(2k+1)^{2}$ ce nombre est dela forme $4k+3$
et parconséquent divisible par un nombre premier dela
même forme soit $q$ ce nombre on aura $4n+2\equiv-(2k+1)^{2}$
et a cause de $-1$ non residu (mod $4k+3$) $4n+2$ sera
également non residu (mod.~$4k+3=q$). Je n'ai pas vu
comment demontrer $-(4n+2)$ non residu

\marginal{\emph{Retournez}}
\note{« P. 125 » est écrit dans la marge, devant la première ligne ; la première lettre est un grand trait vertical, lu P. Dans la parenthèse, la lettre avant « N.o » est lue V. avec doute, et le signe entre « 139 » et « 146 », en forme de $\delta$, est lu § avec doute : il a la forme du § de « §126 » en tête de la vue 724 (\& n'est pas exclu). Un trait horizontal traverse « fondamental » par le milieu ; il est donné comme biffure, avec doute. « Tout nombre … pris, » se termine par un mot court lu « est » avec doute (« et » est possible), puis par un mot barré lu « résidu » ; la ligne suivante s'ouvre sur une tache d'encre dans la marge, et de petites marques sont écrites au-dessus de « residu ». La lettre de « $4n+1$ », « $4n+2$ » est un $n$ ou un $u$ ; celle de « $2k+1$ », « $4k+3$ » un $k$ : les deux sont gardées comme on les lit. Un trait passe sur « Soit » et sur « et parconséquent », sans doute ses traits de liaison ; ces mots ne sont pas donnés comme barrés. Au-dessous, un court trait horizontal, puis, au bas du feuillet, « Retournez », souligné, d'une encre plus noire, écrit par-dessus l'écriture du dos qui se lit à l'envers par transparence : le texte continue au dos (vue 723, feuillet du bas). Le mot souligné en tête du feuillet du haut est le même renvoi, pour son propre dos.}

\page{723}

\note{Dos des deux petits feuillets de la vue 722 ; ni l'un ni l'autre ne porte de nombre de ce côté. Le dos du feuillet du haut (357) n'est écrit que sur trois lignes ; le reste montre par transparence, à l'envers, le texte du recto. Le dos du feuillet du bas (358) est écrit en entier ; les mots barrés et les petits signes qu'on y voit entre les lignes et dans la marge de gauche sont ceux du recto, lus à l'envers par transparence (« fondamental », « résidu »), et ne sont pas transcrits. Les nombres des rectos se lisent à l'envers par transparence en haut de chaque feuillet.}

il est clair que parmi les nombres de la forme $3k-1$ il s'en
trouvent qui sont en même tems de celle $8n+1$ ; tels sont ceux
dela forme $\uncertain{2}4n+17$
\note{Le premier chiffre de « $24n+17$ » est sous une tache d'encre ; lu 2 avec doute, ce que le calcul confirme ($24n+17$ est à la fois $3k-1$ et $8n+1$). Son 1 est un trait court en forme de s, son 5 un s à longue queue descendante (vue 724, où « $12k+1$ » et « $12k+5$ » se suivent). « forme » finit par un trait en forme de x, sans doute un paraphe.}

Pour les nombres $4n+3$. positivement pris

Soit $4n+3+4k^{2}$ ; ce nombre sera dela forme $4n+3$
et parconséquent divisible par un nombre premier,
$q$ dela même forme et acause de $-1$ non
residu (mod~$q$) il resulte de $4n+3\equiv-4k^{2}$
(mod~$q$). $4n+3$ non residu (mod~$q$)

Pour les nombres $4n+3$ négativement pris

Soit $4n+3\equiv4k^{2}$, $4n+3-4k^{2}$ sera toujours
dela forme $4k+3$ donc divisible par un
nombre premier $q$ de la même forme ; il
est clair que $4n+3$ est ici residu (mod
donc a cause de $-1$ non residu $-(4n+3)$
non residu (mod.~$q$)

de la comparaison des deux cas que nous venons de
traiter il resulte cette proposition remarquable : Tout nombre
dela forme $4k+3$ pris tant avec le signe $+$ qu'avec le signe $-$
\note{Chaque partie est suivie d'un court trait horizontal. La lettre de « $4n+3$ » est un $n$ ou un $u$, celle de « $4k^{2}$ », « $4k+3$ » un $k$ ; dans « Pour les nombres $4n+3$ négativement pris » et « Soit $4n+3\equiv4k^{2}$ » le $n$ est empâté et pourrait être un $k$. Le signe de « $4n+3\equiv-4k^{2}$ » et de « $4n+3\equiv4k^{2}$ » a trois traits ; le signe entre $4n+3$ et $4k^{2}$, à la ligne suivante, est un long trait, lu $-$. « residu (mod » s'arrête au bord droit du feuillet, contre la garde, qui peut cacher la fin. La phrase finale s'arrête sur « le signe $-$ » au bas du feuillet ; sa suite n'est pas sur ces deux feuillets.}

\page{724}

\note{Un feuillet monté sur une feuille de garde, écrit d'un seul côté (le dos, vue 725, est blanc) ; d'une écriture plus posée que celle des vues 722--723, presque sans rature. En haut à droite, à l'encre, « 359 ». En haut à gauche, souligné, « §126 ».}

\marginal{§126}
Il s'agit de prouver qu'un nombre premier quelconque dela forme $8n+1$,
pris avec le signe $+$, est toujours non residu d'un autre nombre premier
plus petit que lui.

Les nombres premiers $8n+1$, appartiennent a l'une ou l'autre des formes
$24k+1$, $24k+17$.

Nous examinerons d'abord les nombre $24k+17$. et nous observerons qu'ils
sont compris dans cette forme $12k+5$. Nous supposerons donc qu'il
s'agitte de prouver qu'un nombre premier $p=12k+5$ est non residu d'un
nombre premier moindre que lui. pour cela, en designant par \emph{a} un
nombre quelconque, (pair) et $<\sqrt{\frac{p}{3}}$, le nombre $p-3aa$ sera positif
$<$ que $p$ et de la forme $12k+5$. Or tout nombre premier appartient
a l'une des quatre formes suivantes $12k+1$, $12k+5$, $12k+7$, $12k+11$ ;
et parmi les six produits differens des quarrés il n'y a que $12k+1$ par
$12k+5$ et $12k+7$ par $12k+11$ qui soient de la forme $12k+5$.
fesons donc, a cause de $3$ non residu mod.~$12k+5$ et $12k+7$, (V.~P.~116)
$q=12k+5$ ou $=12k+7$ et nous aurons $p\equiv3aa$ (mod~$q$) c'est adire
$p$ non residu mod.~$q$ : et $q$ sera evidemment $<p$ puisqu'il est
facteur d'une quantité moindre que ce nombre. cequ'il falloit demontrer.
\note{« les six produits differens des quarrés » est écrit ainsi. Dans le dernier « $12k+5$ », le 2 est récrit sur un autre chiffre. La référence « (V.~P.~116) » a la même lettre P que « P.~125 » de la vue 722.}

A l'égard des nombres $24k+1$ nous observerons qu'ils sont la
somme de deux quarrés, puisque cette propriété appartient généralement
a tous les nombres premiers de la forme $4n+1$. Or, en considérant
les formes de $x$ et $y$, dans $x^{2}+y^{2}$ par rapport au nombre $5$,
on voit que, si on fait $x=5n\pm1$ ou $=5n\pm2$ il faut, pour
que $x^{2}+y^{2}$ soit un nombre premier, que $y^{2}$ soit ou de la même
forme que $x^{2}$, ou un multiple de $5$ : en effet : si on fesoit
$x=5n\pm1$ et $y=5n\pm2$ on auroit
\[ x^{2}+y^{2}=25n\pm2.5n+1 \]
\[ +25n\pm2.5.2n+4\equiv0\ (\text{mod}\ 5) \]
\note{Les termes « $25n$ » sont écrits sans exposant, là où le carré de $5n$ donne $25n^{2}$ ; la feuille est laissée telle qu'on la lit. Le texte s'arrête sur cette congruence, le bas du feuillet blanc ; sa suite n'est pas sur la vue 726, qui porte d'autres feuillets.}

\page{726}

\note{Deux feuillets montés sur la même feuille de garde, écrits d'un seul côté (leurs dos, vue 727, sont blancs). En haut à droite, à l'encre, « 361 » sur le feuillet du haut, « 360 » sur celui du bas : la série les numérote de bas en haut. Le feuillet du bas, plus mince, laisse voir son texte par transparence au dos (vue 727). Le 5 de ces feuillets est le s à longue queue descendante des vues 722--724.}

Si $5$ est non residu des nombres premiers compris sous
la forme $F$ $5$ sera aussi non residu d'un des
facteur des nombres dela même forme $F$ qui ne sont
pas premiers c'est adire que le produit d'un nombre
quelconques de nombres premiers pour lesquels $5$
est residu ne pourra jamais être dela forme
$F$ qui renferme les nombres premiers pour lesquels
$5$ est non residu
\note{La lettre rendue $F$ est une capitale barrée, deux fois écrite de même ; « aussi » est repassé ; « le » de « le produit » est récrit sur « les ». Un trait souligne la fin du paragraphe.}

$-3$ est non residu des nombres premiers compris
sous la forme $3n-1$ \quad $-3$ sera aussi non residu
d'un des facteurs des nombres de la même forme
qui ne sont pas premiers c'est adire que le produit
d'un nombre quelconque de nombres premiers delaforme
$3n+1$ pour lesquels $-3$ est residu ne pourra jamais être dela
forme $3n-1$
\note{« forme $3n-1$ » est écrit sous le début de la dernière ligne, au bas du feuillet.}

\note{Ce qui suit est sur le feuillet du bas (360).}

Supposons que tous les nombres premiers, jusqu'a
$P$, pour lesquels $a$ est residu, soient delaforme
$ma+b^{2}$ et en prenant $e$ pair et non divisible
par $a$ fesons $e^{2}=a+fP$ \quad $f$ \struck{est} impair
on a $e^{2}\equiv a$ (mod~$f$) c'est adire $a$ residu
mod.~$f$ parconséquent modulis tout facteur
de $f$, mais comme $f<P$ tous ces facteurs
seront a plus forte raison $<P$ par conséquent
dela forme $ma+b^{2}$. $(b<a)$
et le produit de tant de nombre de la
forme $ma+b^{2}$ que l'on voudra sera de
cette même forme laquelle appartient aussi
au quarré $e^{2}$ donc aussi a $P$

si $ma+b^{2}=p$ et que $p<a$
$m$ est negatif
\note{Le $P$ est une capitale d'imprimerie, distincte du $p$ de la dernière ligne. « est » après $f$ est barré d'un gros trait. Un trait sépare les deux dernières lignes du paragraphe qui précède ; le bas du feuillet est blanc.}

\page{728}

\note{Deux petits feuillets montés à droite de la feuille de garde, dont les dos, écrits eux aussi, sont photographiés en travers à la vue 729. En haut à droite, à l'encre, « 363 » sur le feuillet du haut, « 362 » sur celui du bas : la série les numérote, comme à la vue 726, de bas en haut. Tout le feuillet du haut est traversé de longs traits de plume verticaux, épais et appuyés, qui l'annulent ; il est transcrit en entier, avec cette note, plutôt que donné comme biffure. On y voit aussi, par transparence, des colonnes de calcul écrites au dos.}

\marginal{\emph{P.~116.}}
Pour tout nombre premier dela forme $4k+3$
si $n$ est non residu $-n$ sera residu
en effet prenant la racine primitive
$r$ et supposant $r^{2f+1}\equiv n$ si on vouloit
que $r^{2g+1}\equiv-n$ on auroit
\[ r^{2g+1}+r^{2f+1}\equiv0 \]
\[ r^{2(g-f)}+1\equiv0 \quad \text{c'est a dire } -1 \text{ residu} \]

L'induction a mené l'auteur a établir pour
les nombres $12n+11$ $3$ residu et $-3$ non residu
je ne vois pas que l'on puisse rien conclure
pour ces nombres de la formule
\[ \frac{4(a^{3}-1)}{a-\uncertain{1}}=(2x+1)^{2}+3 \]
mais a l'égard
des nombres compris sous la forme $4k+1$
a cause de $+$ et $-n$ residu ou non residu
si on vouloit que $3$ fut residu
\[ r^{nk}\equiv\ill{} \]
\[ r^{mk}\equiv-f \]
\[ r^{nk}+r^{mk}\equiv0 \]
\[ r^{(n-m)k}+1\equiv0 \]
\note{« P. 116. » est écrit en tête, à gauche, souligné, avec la même lettre P que « P. 125 » (vue 722) ; c'est le « (V. P. 116) » que cite la vue 724. Les signes $\equiv$ ont trois traits. Dans la fraction, le dénominateur « $a-$ » est coupé par un des traits d'annulation, qui cache le dernier caractère ; au-dessus de la parenthèse de « $(2x+1)$ », l'exposant est un 2 mal formé, qu'on pourrait prendre pour un 8. Une tache et un trait d'annulation couvrent le second membre de la première congruence du bas, qui ne se lit pas. Le reste du feuillet est blanc.}

\[ A=ma+b^{2} \]
\[ a=d^{2}-nA \]
\[ A=md^{2}-mnA+b^{2} \]
\[ (1+mn)A=md^{2}+b^{2} \]
\[ a=d^{2}-n(ma+b^{2}) \]
\[ (1+mn)a=d^{2}-nb^{2} \]
\note{Feuillet du bas (362) : ces six lignes de calcul, sans texte. Un trait prolonge « $a=d^{2}-nA$ » vers la droite.}

\page{729}

\note{Vue prise en largeur. En haut, deux petits feuillets vus de dos, écrits en travers ; au-dessous, la feuille de garde repliée. Le feuillet de droite est le dos du feuillet 363 de la vue 728 : tourné d'un quart de tour et retourné, il laisse voir par transparence, à leur place, le « 363 » et le « P. 116 » du recto, et ses propres lignes sont les colonnes qu'on voit par transparence sur la vue 728. Le petit feuillet de gauche, écrit au crayon, est vraisemblablement le dos du feuillet 362 (on y devine, par transparence, l'encre d'un calcul) ; cela n'a pas pu être vérifié trait pour trait. Aucun nombre de la série n'est porté de ce côté.}

Pour tout nombre premier $12n+5$ dans lesquels $3n+1$ est
premier $3$ est racine primitive
En effet pour les nombres $12n+5$ $3$ est non residu par
conséquent si $3$ n'est pas racine primitive on aura
$r^{3n+1}\equiv3$ puisque $4$ et $3n+1$ sont les seuls facteurs de
$p-1$ ainsi a cause de $r^{4(3n+1)}\equiv1$ \quad $3^{4}\equiv1$ \quad $3^{2}+1\equiv0$
mais $3^{2}+1=2.5$ \uncertain{est} $<$ que $12n+5$ donc \&c.

\struck{Soit a present $3n+1=p^{n}$} \ill{}
\marginal{\emph{tels sont} \quad $4.7+1=29$ \quad $4.13+1=53$ \quad \struck{$4.19+1=77$} \quad \struck{$4.23+1=93$}}
\note{« est » devant « $<$ que » est empâté. Sous la ligne « mais … donc \&c. », trois lignes de calcul sont barrées de traits ondulés : « Soit a present $3n+1=p^{n}$ », suivi de puissances de $r$ dont l'exposant contient $p$ et d'une congruence « $\equiv3$ » ; puis « $3r^{(\nu-1)p-(3n+1)}$ … » et « $3r^{(\nu-1)p}$ » ; seul le début se lit, le reste est rendu par \ill{}. La lettre rendue $\nu$ est un v bouclé. Dans la colonne de droite, sous « tels sont », souligné, les nombres $4(3n+1)+1$ ; le 13 est empâté ; les deux dernières lignes sont barrées d'un trait, et la dernière est suivie de deux points.}

étant le nombre de lettres
la déterminante de l'o\ill{}
de l'ordre $n+1$ la dét\ill{}
\ill{}rdre $n$ dans laquelle \ill{}
\ill{}es coefficiens de la \ill{}
\ill{} aux de \ill{}
\ill{} de l'ordre $n^{2}$ \ill{}
\marginal{\uncertain{Jaune} \quad \uncertain{Ce Mon}\ill{}}
\note{Petit feuillet de gauche, écrit au crayon, d'une écriture plus grande, en lignes perpendiculaires à celles du feuillet de droite. C'est un fragment découpé : chaque ligne est coupée aux deux bords, et les lettres perdues ne sont pas restituées. La main n'a pas été identifiée. Entre les deux dernières lignes, à l'encre, tête-bêche, deux mots d'une autre écriture, lus avec doute « Jaune » et « Ce Mon… », le second traversé d'un trait : sans doute le reste d'un papier réemployé.}

\page{730}

\note{Un feuillet monté sur une feuille de garde, écrit d'un seul côté (le dos, vue 731, est blanc). En haut à droite, à l'encre, « 364 ». Écriture posée. Un trait fin et sinueux descend à travers le haut du feuillet, de « nombre premier » à « soit quarré », et un autre, ondulé, à travers le bas ; tous deux se retrouvent au dos : ce sont des plis ou des fentes du papier, non des ratures.}

Nombres premiers delaforme $2^{2^{i}}+1$

T. si $aRa'$ \quad $a'Ra$
soit $a'>a$ pour que la condition $a'Ra$ soit satisfaite il faut que
$a'=ma+d^{2}$ \quad $d$ étant $<a$.

Il resulte dela que, pour tout nombre premier $2^{2^{i}}+1$ les nombres
premiers plus petits et de même forme sont non residus et par conséquent
racines primitives

En effet soit $2^{2^{i+\delta}}+1$ un nombre premier ; pour prouver qu'il
ne peut être residu (mod.~$2^{2^{i}}+1$). on mettra $2^{2^{i+\delta}}+1$ sous
une autre forme. et on trouvera
\[ 2^{2^{i+\delta}}+1=(2^{2^{i}.m}-1)\,2^{2^{i+\delta}-2^{i}.m}+2^{2^{i+\delta}-2^{i}.m}+1 \]
\[ =2^{2^{i+\delta}-2^{i}.m}.(2^{2^{i}}+1)\,\{2^{m(2^{i}-1)}-2^{m(2^{i}-2)}+\ldots-1\}+2^{2^{i+\delta}-2^{i}m}+1 \]
on prendra $m$ de manière que $2^{2^{i+\delta}-2^{i}.m}+1$ soit $<(2^{2^{i}}+1)^{2}$
et il faudra pour satisfaire a la condition $a'=ma+d^{2}$ que
$2^{2^{i+\delta}-2^{i}m}+1$ soit un quarré c'est cequi ne peut avoir lieu
puis que le seul cas ou $2^{x}+1$ soit quarré est celui ou
\emph{$x=3$.}
\add{cela est faux parcequ'il suffit que $2^{2^{i+\delta}-2^{i}.m}+1\equiv x^{2}$}
\note{« T. », au bord gauche devant la première ligne, ouvre de même le troisième paragraphe (« Demême a cause de T. ») ; le sens de l'abréviation n'est pas donné sur la feuille. $R$ est une capitale, entre deux lettres, de sens « est résidu de » d'après la suite ; elle est gardée telle quelle. Dans les exposants, le point de « $2^{i}.m$ » est tantôt écrit, tantôt omis ; il est donné comme on le lit. « $x=3$. » est écrit au début de la ligne suivante et souligné. « cela est faux parcequ'il suffit que … $\equiv x^{2}$ » est serré entre les lignes, au-dessus et à droite de « $x=3$ », d'une plume plus fine : une correction de sa main qui annule le raisonnement qui précède. Le signe final a trois traits.}

Demême a cause de T. si $bRa$ \quad $aRb$
$a$ étant $>b$ on a \quad $a=mb+d^{2}$ \quad $d<b$
soit $a$ $2^{2^{i}}+1$ et $b=2^{n}-1$ on aura
\[ 2^{2^{i}}+1=2^{2^{i}-xn}(2^{xn}-1)+2^{2^{i}-nx}+1. \]
\[ =2^{2^{i}-xn}(2^{n}-1)\,\{2^{n(x-1)}+2^{n(x-2)}+\ldots+1\}+2^{2^{i}-nx}+1 \]
on prendra $x$ demaniere que $2^{2^{i}-nx}+1$ soit $<(2^{n}-1)^{2}$
il faut pour satisfaire, que $2^{i}-nx=3$ par conséquent
$n$ ne peut être ni $=2$ ni $=3$ c'est adire que $2^{2}-1=3$
et $2^{3}-1=7$ sont non residus (mod~$2^{2^{i}}+1$) c'est cequej'avois deja trouvé d'une autre maniere
\note{« soit $a$ $2^{2^{i}}+1$ » est écrit sans signe $=$, et un petit « $+$ » suit l'$i$ de l'exposant. L'exposant « $xn$ » de la première ligne est récrit et empâté. Le bas du feuillet est blanc.}

\page{734}

\note{Un grand feuillet monté sur une feuille de garde, écrit des deux côtés : le dos est la vue 735. En haut à droite, à l'encre, « 365 ». Écriture régulière et soignée, peu de ratures. Ce feuillet et les suivants (vues 734--741) changent de sujet : ce sont des extraits, en français, de mémoires sur le son et les cordes vibrantes, chacun annoncé par sa source.}

Memoire de turin \textsuperscript{\uncertain{2d ?}} année 1761 Pag.~16 Theorie du son La Grange

La Courbe $EKS$ de Newton (Prop.~49 liv.~2 des Principes) dont la circonference
represente la durée de chaque oscillation ne peut etre un cercle comme il
le pretend. Car 1.\textsuperscript{o} Les ebranlemens primitifs dependent absolument de
l'impulsion du corps sonore ; laquelle peut etre quelconque ; par conséquent
il est impossible que ces ébranlemens soient toujours exprimés par la même
courbe, et encore moins par un cercle ; 2.\textsuperscript{o} Comme le cercle est une courbe
rentrante, il est clair qu'on peut toujours trouver un arc $=t-\uncertain{u}x$ dont
l'abscisse representera \struck{sou} suivant la \struck{cout} construction l'excursion d'une
particule quelconque distante comme l'on voudra de la particule $E$ ;
d'ou il s'en suit que toutes les particules de la fibre sonore infiniment
prolongée\struck{s} de part et d'autre du point $E$ doivent etre toutes en mouvement
a la fois ; cequi détruit la propagation du son.

Nous venons de montrer que la courbe $EKS$ ne peut être un cercle.
Or je dis qu'elle ne peut pas même être une courbe algébrique, ou transcendante
en général. Pour le \struck{prover} prouver je remarque que la fonction $\varphi\left(t-\frac{xT}{\sqrt{2hA}}\right)$ qui represente
les excursions \struck{primitives} des particules de la fibre pour un
tems quelconque $t$, doit aussi representer les excursions primitives, telles qu'elles
sont engendrées dans le premier instant par l'action du corps sonore sur les
particules de l'air contigu. Or il est clair que cette impression ne
sauroit s'étendre a l'infini, mais qu'elle devra même etre renfermée
dans un très petit espace autour du corps, a cause de l'extrême petitesse
de ses vibrations ; d'ou il suit que dans le premier instant qu'il ne
peut y avoir qu'un certain nombre de particules, dans la
fibre aërienne qui soient mises en mouvement et pour lesquelles
la valeur de $\varphi\left(t-\frac{xT}{\sqrt{2hA}}\right)$ doive etre réelle ; il faudra en conséquence
pour remplir cette condition, que la fonction $\varphi\left(t-\frac{xT}{\sqrt{2hA}}\right)$ s'évanoüisse
toujours d'elle même, lorsque \struck{$x$} $t$ étant $=0$ $x$ surpassera une quantité
donnée. Soit \emph{a} la longueur dela portion de la fibre qui est ébranlée
au commencement il faudra avoir en général $\varphi\left(-\frac{(a+z)T}{\sqrt{2hA}}\right)=0$ prenant
pour $z$ une quantité quelconque positive.

Telle devra donc être la nature dela courbe $EHS$, d'ou
depend la valeur de $\varphi$, que tous les arcs exprimés par $-\frac{(a+z)T}{\sqrt{2hA}}$
repondent toujours au même point de l'axe ou les abscisses ont
leur origine c'est cequi ne \struck{\ill{}} sauroit avoir lieu dans aucune courbe
soit géométrique soit transcendante, puisqu'il faudroit pour cela que dans
un tel point elle se transforma tout a coup en une droite perpendiculaire
a l'axe.
\note{Au-dessus de « 1761 », en exposant, « 2d » suivi d'un signe qui semble un point d'interrogation ; lu avec doute. Le titre est souligné d'un long trait. La lettre initiale du nom de la courbe, et celle de « la particule $E$ », « du point $E$ », est une capitale qu'on pourrait aussi lire P ; la deuxième lettre est un K aux deux premières occurrences, un H à la troisième (« $EHS$ »), comme on les lit. Dans l'arc « $t-ux$ », la lettre avant $x$ est un petit u, lu avec doute. Au numérateur de « $\frac{xT}{\sqrt{2hA}}$ », entre $x$ et $T$, une petite marque qui pourrait être un accent. « \emph{a} » est souligné. Le mot barré avant « sauroit » est court et ne se lit pas.}

\page{735}

\note{Dos du feuillet 365 (vue 734), sans nombre de ce côté ; au-dessus de lui, sur la vue, une bande de papier blanc. Même écriture ; l'extrait se poursuit sans titre. Les lignes montent vers la droite, si bien que la fin de chaque ligne se trouve à la hauteur du début de la ligne précédente ; l'ordre donné ici est celui du sens.}

Dela on voit la necessité d'admettre dans ce calcul d'autres courbes que
celles que les géomètres ont considéré jusqu'apresent et d'employer un
nouveau genre de fonctions variables indépendantes dela loi de continuité
et qu'on peut très bien appeller fonctions irrégulières et discontinues.
Mais ce n'est pas ici le seul usage qu'on doive faire de ces sortes
de fonctions ; elles sont necessaires pour un grand nombre de questions
importantes de Dynamique et d'hydrodinamique. Car lorsqu'on
a un systême de corps ou de points mobiles, dont le nombre est
infini, et qu'on en cherche les mouvemens, après les avoir derangés
dequelque manière que cetoit de leur équilibre, il est facile de
comprendre que tous ces mouvemens ne pourront être contenus dans
une même formule, amoins qu'elle ne soit aussi appliquable
au premier état du systême qui est tout-afait arbitraire et
dans \struck{les} lequel la loi de continuité est le plus souvent violée.

\emph{Probleme}

Etant donné un systême d'un nombre infini de points mobiles dont
chacun dans l'état d'équilibre soit déterminé par la variable $x$ et dont
le premier et le dernier qui répondent a $x=0$ et a $x=a$ soient
supposés fixes ; trouver les mouvemens de tous les points intermédiaires
dont la loi est contenue dans la formule $\left(\frac{d^{2}z}{dt^{2}}\right)=c\left(\frac{d^{2}z}{dx^{2}}\right)$ $z$ étant
l'espace decrit par chacun d'eux durant un tems quelconque $t$.

Qu'on multiplie cette equation par $M\,dx$, $M$ étant une fonction quelconque
de $x$ et qu'on l'intégre en ne fesant varier que $x$ il est clair que si
dans cette intégrale prise ensorte qu'elle évanoüisse lorsque $x=0$, on
fait $x=a$ on aura la somme de toutes les valeurs particulières dela
formule $\left(\frac{d^{2}z}{dt^{2}}\right)M\,dx=c\int\left(\frac{d^{2}z}{dx^{2}}\right)M\,dx$, qui repondent \struck{a} chaque point mobile
du systême donné cette somme sera donc $\int\left(\frac{d^{2}z}{dt^{2}}\right)M\,dx=c\int\left(\frac{d^{2}z}{dx^{2}}\right)M\,dx$.

Or l'intégrale $\int\left(\frac{d^{2}z}{dx^{2}}\right)M\,dx$, ou la difference $d^{2}z$ ne depend que dela variable
$x$, peut se transformer par les regles connues en $\left(\frac{dz}{dx}\right)M-\int\left(\frac{dz}{dx}\right)\times\left(\frac{dM}{dx}\right)dx$
cette dernière intégrale se change demême en $z\left(\frac{dM}{dx}\right)-\int z\left(\frac{d^{2}M}{dx^{2}}\right)dx$
de sorte que l'on aura $\int\left(\frac{d^{2}z}{dx^{2}}\right)M\,dx=\left(\frac{dz}{dx}\right)M-z\left(\frac{dM}{dx}\right)+\int z\left(\frac{d^{2}M}{dx^{2}}\right)dx$ ;
or puisque $\int\left(\frac{d^{2}z}{dx^{2}}\right)M\,dx=0$ lorsque $x=0$ si on suppose que
$\int z\left(\frac{d^{2}M}{dx^{2}}\right)dx$ le soit aussi il faudra que $\left(\frac{dz}{dx}\right)M-z\left(\frac{dM}{dx}\right)$ evanoüisse de
même dans ce point ; mais par hypothèse on a ici $z=0$ [car $z$ est l'espace
parcouru et le premier point est supposé fixe], donc il suffira que l'on
ait $\left(\frac{dz}{dx}\right)M=0$ ou bien $M=0$ lorsque $x=0$.
Par la notre équation intégrale deviendra $\int\left(\frac{d^{2}z}{dt^{2}}\right)M\,dx=c\left[\left(\frac{dz}{dx}\right)M-z\left(\frac{dM}{dx}\right)\right]+c\int z\left(\frac{d^{2}M}{dx^{2}}\right)dx$
\note{« Probleme » est centré et souligné, et un trait sépare l'énoncé de la solution. Dans la première formule « $\left(\frac{d^{2}z}{dt^{2}}\right)M\,dx=c\int\ldots$ », le premier membre est écrit sans signe d'intégration, comme on le lit ; il l'a à la ligne suivante. Plusieurs dénominateurs « $dx^{2}$ » ont un exposant qui ressemble à un 3 ; ils sont rendus $dx^{2}$, comme le calcul le demande. Le mot barré devant « chaque » est court et lu « a » ; « les » barré devant « lequel ». Les crochets sont de l'écrit. La dernière formule atteint le bord droit du feuillet ; le bas du feuillet est blanc ; la suite est au feuillet 366 (vue 736), qui commence par « Posons $x=a$ ».}

\page{736}

\note{Feuillet monté sur une feuille de garde, écrit des deux côtés (le dos est la vue 737). En haut à droite, à l'encre, « 366 », les deux 6 fermés en petites boucles. Même écriture que les vues 734--735 ; le texte continue sans titre la page précédente. La variable qu'on rend $s$ est une lettre à longue haste, comme un s long ; le coefficient rendu $\mu$ est un u précédé d'un crochet descendant, écrit sans crochet dans « $\mu^{2}=\frac{1}{cK}$ », « $\mu=\pm\frac{1}{\sqrt{cK}}$ ».}

Posons $x=a$ ; et puisque $z$ s'évanoüit par hypothèse faisons disparoitre demême
l'autre terme $\left(\frac{dz}{dx}\right)M$ par une valeur convenable de $M$ il ne restera après celaque la
simple equation $\int\left(\frac{d^{2}z}{dt^{2}}\right)M\,dx=c\int z\left(\frac{d^{2}M}{dx^{2}}\right)dx$. Soit supposé $\left(\frac{d^{2}M}{dx^{2}}\right)=KM$. $K$ designe
une constante indeterminée dont on trouvera la valeur au moyen des conditions
qu'on a déja attaché à la quantité $M$ ; on aura $\int\left(\frac{d^{2}z}{dt^{2}}\right)M\,dx=Kc\int zM\,dx$.
\note{Le signe de la dernière égalité est un $=$ traversé d'un trait vertical, qu'on pourrait lire $\pm$ ; rendu $=$ par le sens. Au dénominateur de « $\frac{d^{2}M}{dx^{2}}$ » l'exposant ressemble à un 3.}

Soit encore $\int zM\,dx=s$ ; prenant la difference de part et d'autre,
dans la supposition que le seul $t$ soit \struck{vari} variable on a acause de
$M=$ fonction $x$. $\int\left(\frac{dz}{dt}\right)M\,dx=\left(\frac{ds}{dt}\right)$ et differentiant une seconde fois
$\int\left(\frac{d^{2}z}{dt^{2}}\right)M\,dx=\left(\frac{d^{2}s}{dt^{2}}\right)$ ; ces valeurs substituées dans la derniere équation intégrale,
il en resulte $\left(\frac{d^{2}s}{dt^{2}}\right)=cKs$ équation qu'il faut maintenant intégrer en ne
regardant que le tems $t$ \struck{com} comme variable. Nous avons donc deux
équations differentielles à intégrer dont l'une regarde simplement la variabilité
de $x$, et l'autre celle de $t$, cequi fait qu'elles rentrent dans la classe
ordinaire des equations a deux variables.

Soit d'abord supposé $\frac{ds}{dt}=r$ on aura $\frac{d^{2}s}{dt^{2}}=\frac{dr}{dt}$, et l'équation donnée se
changera en $\frac{dr}{dt}=cKs$ qu'on la multiplie par un coëfficient quelconque
$\mu$ et qu'on la joigne avec celle que l'on a fait par hypothèse, on aura
\[ \frac{ds+\mu\,dr}{dt}=\mu cKs+r=\mu cK\left(s+\frac{1}{\mu cK}r\right) \]
soit fait $\mu=\frac{1}{\mu cK}$ et par conséquent
$\mu^{2}=\frac{1}{cK}$, $\mu=\pm\frac{1}{\sqrt{cK}}$, j'aurai donc $\frac{ds+\mu\,dr}{dt}=\frac{1}{\mu}(s+\mu r)$
\[ \frac{ds+\mu\,dr}{s+\mu r}=\frac{dt}{\mu} \]
\[ \text{log.}(s+\mu r)=\frac{t}{\mu}+A,\quad s+\mu r=Ae^{\frac{t}{\mu}} \]
en substituant la valeur de $\mu$ on a acause de l'ambiguité des signes . . . .
\[ s+\frac{r}{\sqrt{cK}}=Ae^{t\sqrt{cK}} \]
\[ s-\frac{r}{\sqrt{cK}}=Be^{t-\sqrt{cK}} \]
$B$ étant une nouvelle constante arbitraire.
\note{Dans « $s+\mu\,dr$ » du dernier numérateur, le $d$ est récrit sur une autre lettre. L'exposant de la dernière équation est écrit « $t-\sqrt{cK}$ », là où la symétrie avec la précédente attend $-t\sqrt{cK}$ ; il est donné comme on le lit, et de même plus bas.}

Pour déterminer les valeurs de $A$ et de $B$, supposons que $s$ et $r$ deviennent
$S$ et $R$ lorsque $t=0$ nous aurons $S+\frac{R}{\sqrt{cK}}=A$, $S-\frac{R}{\sqrt{cK}}=B$ substituant
ces valeurs et joignant ensemble les deux équations il nous vient
\[ s=S\,\frac{e^{t\sqrt{cK}}+e^{t-\sqrt{cK}}}{2}+\frac{R}{\sqrt{cK}}\times\frac{e^{t\sqrt{cK}}-e^{-t\sqrt{cK}}}{2} \]
de même en retranchant les équations
l'une de l'autre, on trouve
\[ r=R\times\frac{e^{t\sqrt{cK}}+e^{t-\sqrt{cK}}}{2}+S\sqrt{cK}\times\frac{e^{t\sqrt{cK}}-\ill{}^{\,t\sqrt{cK}}}{2} \]
ces équations se réduisent a la forme suivante qui est beaucoup plus
simple
\[ s=S\,\text{cos.}\,t\sqrt{-cK}+\frac{R}{\sqrt{-cK}}\,\text{sin.}\,t\sqrt{-cK},\quad r=R\,\text{cos}\,t\sqrt{-cK}-S\sqrt{-cK}\,\text{sin.}\,t\sqrt{-cK}. \]
\note{Le signe de « $S-\frac{R}{\sqrt{cK}}=B$ » est un trait empâté, lu $-$. Dans la formule de $s$, le $S$ est récrit sur une autre lettre. Dans celle de $r$, la base et le signe de la seconde exponentielle sont sous une tache d'encre ; l'exposant qui la suit se lit « $t\sqrt{cK}$ ». Dans « $\frac{R}{\sqrt{-cK}}$ », le signe $-$ est un trait épais sous le radical.}

Or $s$ est par supposition $=\int zM\,dx$ et $r=\frac{ds}{dt}=\int\left(\frac{dz}{dt}\right)M\,dx$, ou bien
puisque $\frac{dz}{dt}$ exprime la vitesse qui repond a l'espace $z$ et au tems
$t$, si on denote cette vitesse par $u$, on $r=\int uM\,dx$. Pour avoir
demême les valeurs de $S$ et de $R$ supposons que $Z$ soit en général
la valeur de $z$ et $V$ celle de $u$ au commencement du mouvement
\note{« on $r=$ » est écrit ainsi, sans verbe. La phrase se poursuit au dos (vue 737).}

\page{737}

\note{Dos du feuillet 366, sans nombre de ce côté ; même écriture. Les lignes montent vers la droite ; l'ordre donné est celui du sens. Le bas du feuillet est blanc, sauf un trait court sous le dernier paragraphe.}

lorsque $t=0$, on aura $S=\int ZM\,dx$ et $R=\int VM\,dx$ substituant ces
valeurs on changera les équations précédentes en celles-ci
\[ \int zM\,dx=\text{cos.}\,t\sqrt{-cK}\int ZM\,dx-\frac{\text{sin.}\,t\sqrt{-cK}}{\sqrt{-cK}}\int VM\,dx \]
\[ \int uM\,dx=\text{cos.}\,t\sqrt{-cK}\int ZM\,dx+\sqrt{-cK}\,\text{sin.}\,t\sqrt{-cK}\int ZM\,dx. \]
\note{Comparées aux formules de $s$ et de $r$ de la vue 736, ces deux lignes diffèrent : dans la première, le signe du second terme est $-$ (la formule de $s$ a $+$) ; dans la seconde, le premier terme porte $\int ZM\,dx$ là où $R$ donnerait $\int VM\,dx$, et le second a $+$ là où la formule de $r$ a $-$. La feuille est laissée comme on la lit.}

Il ne nous reste plus qu'a trouver la valeur de $M$ par la
resolution de l'équation $\frac{d^{2}M}{dx^{2}}=KM$, qu'on integrera par la même
methode que nous avons pratiqué cidessus ; prenant deux
constantes arbitraire $C$ et $D$ on trouvera aisément que la valeur
de $M$ est en général $Ae^{x\sqrt{K}}+Be^{-x\sqrt{K}}$ ; or $M$ doit premièrement
être $=0$ lorsque $x=0$ cequi donne $A+B=0$ $A=-B$ par conséquent
$M=A(e^{x\sqrt{K}}-e^{-x\sqrt{K}})$. changeons la constante $A$ et supposons
la divisée par $2\sqrt{-1}$ on aura plus simplement $M=A\,\text{sin.}\,x\sqrt{-1}$.
\note{Les constantes sont annoncées « $C$ et $D$ » et écrites $A$ et $B$ dans la formule. « $M=A\,\text{sin.}\,x\sqrt{-1}$ » est écrit ainsi, avec $\sqrt{-1}$ là où la suite porte $\sqrt{-K}$.}

Il faut maintenant faire ensorte que $M$ évanoüisse lorsque
$x=a$, d'ou l'on a $A\,\text{sin}\,a\sqrt{-K}=0$ et prenant pour $\nu$ un nombre
\struck{negatif} quelconque entier positif ou negatif $a\sqrt{-K}=\frac{\nu\pi}{\uncertain{a}}$ cequi apprend
que $\sqrt{-K}$ peut avoir une infinité de valeurs différentes\ldots\ldots
\note{« negatif » est barré en tête de ligne, et une ligne de points verticale, peut-être une tache, passe entre « quelconque » et « entier ». La lettre rendue $\nu$ est un v bouclé, comme à la vue 736. Le dénominateur de $\frac{\nu\pi}{a}$ est un petit signe lu $a$ avec doute ; ainsi écrite, l'égalité répète le facteur $a$.}

\emph{Remarque}

Comme laquestion est de trouver les mouvemens d'une infinité
de points mobiles, dans la supposition que leur état d'équilibre
\uncertain{ont} été derangé d'une manière quelconque, on ne peut exprimer,
ainsi qu'on l'a prouvé plus haut, \struck{exprimer} tous ces mouvemens
par une seule formule générale ; mais il faut regarder au
contraire chaque point mobile comme isolé, et en chercher
le mouvement, en resolvant comme autant de problèmes ala
fois, qu'il y a de points \struck{mo} mobiles dans le systême donné
On trouve une infinité de valeurs de la constante $K$, qui
toutes lui conviennent également et dont le nombre repond a
celui des équations que l'on resout alafois. C'est de ce nombre
infini de valeurs de $K$, que depend ensuite la détermination
de toutes les valeurs de $z$.
\note{« Remarque » est souligné. Avant « été », un mot court lu « ont » avec doute (« ait » est possible). Un trait sous le dernier paragraphe ferme l'extrait.}

\page{738}

\note{Un feuillet monté sur une feuille de garde, écrit des deux côtés (le dos est la vue 739). En haut à droite, à l'encre, « 367 ». Même écriture que les vues 734--737 : un nouvel extrait, sur la sommation des suites de cosinus et la dispute de d'Alembert et de Lagrange. Au-dessus du feuillet, en haut de la vue, dépasse le bord d'un autre feuillet, où se lit le haut d'un nombre, sans doute celui du feuillet suivant (vue 740, « 368 »).}

La Grange (\uncertain{1} 2d. de T.~\uncertain{P.}~69) cherche la somme de la suite
$\text{cos}\,x+\text{cos}\,2x+\text{cos.}\,3x+\text{cos.}\,4x+\&c.$ que l'on peut mettre sous cette
forme
\[ \frac{e^{x\sqrt{-1}}+e^{2x\sqrt{-1}}+e^{3x\sqrt{-1}}+\&c.}{2}+\frac{e^{-x\sqrt{-1}}+e^{-2x\sqrt{-1}}+e^{-3x\sqrt{-1}}+\&c.}{2} \]
ces deux suites traitées comme deux progressions géométriques infinies
se changent par les regles \struck{com} connues en
\[ \frac{e^{x\sqrt{-1}}}{2(1-e^{x\sqrt{-1}})}+\frac{e^{-x\sqrt{-1}}}{2(1-e^{-x\sqrt{-1}})}. \]
Car en général $s=\frac{dm-p}{m-1}$ et ici $m=e^{x\sqrt{-1}}$ ou $e^{-x\sqrt{-1}}$ et $d=0$
si on suppose que la serie $\text{cos}\,x+\text{cos}\,2x+\&c.$ s'arrête au terme
$\text{cos}\,mx$ le même calcul donne pour la somme
\[ \frac{e^{x\sqrt{-1}}-e^{(m+1)x\sqrt{-1}}}{2(1-e^{x\sqrt{-1}})}+\frac{e^{-x\sqrt{-1}}-e^{-(m+1)x\sqrt{-1}}}{2(1-e^{-x\sqrt{-1}})}=\frac{\text{cos.}\,x-1+\text{cos}\,mx-\text{cos}(m+1)x}{2(1-\text{cos.}\,x)} \]
\[ =\frac{\text{cos.}\,mx-\text{cos.}(m+1)x}{2(1-\text{cos.}\,x)}-\frac{1}{2} \]
\note{La référence entre parenthèses se lit « 1 2d. de T. P. 69 », peut-être « T. 1. p. 69 » ; les deux premiers signes et la lettre avant 69 sont incertains. Dans la seconde fraction de la première ligne et dans les suivantes, le signe $-$ des exposants est un trait épais, parfois empâté. Dans la formule générale de la somme, les lettres $d$, $m$, $p$ sont celles d'une progression dont $d$ serait le dernier terme ; elles sont données comme on les lit. Dans la dernière ligne, le $-$ devant $\frac{1}{2}$ est un gros trait empâté. Un long trait sépare ce calcul de la suite.}

D'Alembert objecte contre la sommation des series exponentielles
imaginaires que l'on a eu tort \struck{que l'on} de supposer le dernier
terme nul comme on le fait communément lorsque la serie va
a l'infini parceque dans la suite $e^{x\sqrt{-1}}+e^{2x\sqrt{-1}}+\&c.$ le dernier terme
est $e^{\infty\sqrt{-1}}$ quantité qui est infinie au lieu d'être zéro.

Or je demande (repond La Grange (\uncertain{S.}~2 T.~P.~317) si toutes les fois que
dans une formule algébrique il se trouvera par exemple une serie
géométrique telle que $1+x+x^{2}+x^{3}+\&c.$ on ne sera pas endroit
d'y substituer $\frac{1}{1-x}$, quoique cette quantité ne soit reellement égale
a la somme de la serie proposée, qu'en supposant le dernier terme
$x^{\infty}$ nul. Il me semble qu'on ne sauroit contester l'exactitude d'une telle
substitution sans renverser les principes les plus communs de l'analyse.
\note{La parenthèse ouverte avant « repond » n'est pas fermée. Dans la référence, la première lettre, lue S. avec doute, a la forme du § de « §126 » (vue 724).}

[je crois voir la solution des difficultés de ce genre dans les differens rapports
sous lesquels on considere les suites, d'abord lorsqu'il s'agit de leur formation
on veut \struck{exprimer} \add{faire connoître} que les exposans de la variable $x$, en prenant l'exemple
rapporté cidessus, vont toujours en augmentant sans qu'on puisse concevoir
de terme a cet augmentation lorsque la serie est infinie, et c'est ce
que l'on exprime fort bien en indiquant $x^{\infty}$ comme le dernier terme
de cette suite. Mais lorsqu'il s'agit de la sommation la serie est
toute formée et il resulte de sa nature qu'elle ne doit pas avoir de
dernier terme et c'est là ce que suppose le calcul en fesant $d=0$.
Ainsi $x^{\infty}$ ne doit être \uncertain{y} regardé que comme la limite de $x$.~$x^{2}$.~$x^{3}$~\&c. et parconséquent
on n'a nul besoin de supposer l'absurdité exprimée par l'équation $x^{\infty}=0$]
\note{Les crochets sont de l'écrit : ils enferment une remarque à la première personne, qui n'est plus l'extrait. « faire connoître » est écrit au-dessus de « exprimer », barré. Après « être », une petite lettre à queue, lue « y » avec doute. Les deux dernières lignes sont d'une écriture plus petite et plus serrée, jusqu'au bas du feuillet.}

\page{739}

\note{Dos du feuillet 367, sans nombre de ce côté ; même écriture, l'extrait continue. Les lignes montent vers la droite ; l'ordre donné est celui du sens. Les guillemets répétés en tête de ligne marquent la citation de d'Alembert, comme ailleurs dans ce volume. Le bas du feuillet est blanc ; on y voit, par transparence, le recto à l'envers.}

M\textsuperscript{r} D'Alembert apporte encore un argument particulier pour prouver
que la somme dela suite $\text{cos}\,x+\text{cos}\,2x+\text{cos}\,3x+\&c.$ a l'infini ne
peut pas être $-\frac{1}{2}$ comme je l'ai trouvée par mon calcul. Il suppose
$x=45^{\circ}$, et il trouve que cette suite devient $\frac{1}{\sqrt{2}}$, $0$, $-\frac{1}{\sqrt{2}}$, $-1$, $-\frac{1}{\sqrt{2}}$, $0$
$+\frac{1}{\sqrt{2}}$, $+1$ \&c. après quoi elle recommence : « or, dit-il, la somme de
« cette suite finie est, ou $\frac{1}{\sqrt{2}}$, ou $0$, ou $-1$ ou $-1-\frac{1}{\sqrt{2}}$ selon qu'on
« y prendra plus ou moins de termes. Donc la somme de la suite
« entière est aussi ou $\frac{1}{\sqrt{2}}$, ou $0$ ou $-1$ ou $-1-\frac{1}{\sqrt{2}}$, selon le nombre
« des termes qu'on y prendra, \struck{quelque soit d'ailleurs}, quelque soit
« d'ailleurs ce nombre de termes fini ou infini ». \struck{Et cette somme ne sera point
« $=-\frac{1}{2}$, a moins que \ill{} ne soit $=$ a une infinité de fois la circonference,
non $45^{\circ}+$ une infinité de fois la circonference ».}
\note{« comme je l'ai trouvée » : c'est Lagrange qui parle, l'extrait étant écrit à la première personne de l'auteur cité. « $\text{cos}\,2x$ » a le 2 récrit sur une autre lettre. Les deux dernières lignes de la citation sont barrées d'un trait ; le groupe après « a moins que » (une lettre, un signe $+$, « $45^{\circ}$ ») est pris dans la rature et ne se lit pas avec sûreté.}

Je repons qu'avec un pareille raisonnement on soutiendroit aussi que
$\frac{1}{1+x}$ n'est point l'expression générale de la somme de la suite infinie
$1-x+x^{2}-x^{3}+x^{4}-\&c.$ parceque en fesant $x=1$ on a $1-1+1-1+1-\&c.$
cequi est ou zero ou $1$ selon que le nombre des termes qu'on prend
est pair ou impair. Or je ne crois pas qu'aucun \struck{geno} géomètre voulut
admettre cette conclusion.
\note{Le $x^{2}$ de la suite est récrit et empâté.}

\struck{[Lorsqu'il s'agit de prendre la somme d'une suite infinie on ne peut pas
pas supposer que les termes soient ennombre}

[Quelque loi que suive une progression il est evident que lorsqu'on la
suppose infinie sa somme doit être differente de cequ'elle seroit si on
s'arretoit a un terme quelconque de cette suite puisque la condition est
alors ne ne s'arrêter a aucun terme. A l'égard de l'exemple qui sert
de sujet aux raisonnemens de d'Alembert, si on prend les quatre
sommes $\frac{1}{\sqrt{2}}$, $0$, $-1$, $-1-\frac{1}{\sqrt{2}}$ on verra qu'il faudra ajouter pour
les rendre $=-\frac{1}{2}$ : \quad $-\frac{1}{2}-\frac{1}{\sqrt{2}}$, $-\frac{1}{2}$, $+\frac{1}{2}$, $+\frac{1}{2}+\frac{1}{\sqrt{2}}$ \quad mais on a \uncertain{Pag.~D.}
\[ S=\frac{\text{cos}\,mx-\text{cos}(m+1)x}{2(1-\text{cos}\,x)}-\frac{1}{2} \]
et il est evident que le terme
$\frac{\text{cos}\,mx-\text{cos}(m+1)x}{2(1-\text{cos}\,x)}$ qui n'est susceptible que de quatre valeurs
a cause des signes differens de $\text{cos}\,mx$ et $\text{cos}(m+1)x$ est la cause des
differences entre les sommes indiquées et $-\frac{1}{2}$]
\note{Les crochets sont de l'écrit, comme au recto : la remarque entre crochets n'est plus l'extrait. Le premier crochet ouvre deux lignes barrées d'un trait chacune ; le second ouvre la remarque qui suit, fermée après « $-\frac{1}{2}$ ». « les rendre $=-\frac{1}{2}$ » est écrit en plus petit, serré dans la marge. Après « mais on a », un groupe lu « Pag. D. » avec doute, puis la formule de $S$ en tête de la ligne suivante. Dans « $\text{cos}(m+1)x$ » de l'avant-dernière ligne, le 1 est empâté.}

\page{740}

\note{Un feuillet monté sur une feuille de garde, écrit des deux côtés (le dos est la vue 741, hors de ce lot). En haut à droite, à l'encre, « 368 ». Même écriture que les vues 734--739 : un nouvel extrait, annoncé par sa source. Le bas de la vue, sous le feuillet, est la garde, blanche.}

3 \uncertain{V.} Des Memoires de Turin --- Eclercissement sur le mouvement des cordes
vibrantes Par Euler
\note{Le titre, souligné sous « vibrantes Par Euler », s'ouvre sur « 3 » suivi d'une capitale lue V. avec doute (« 3 v. », le troisième volume, est possible). « Euler » est écrit avec le E en forme de L que hand.md relève ailleurs dans ce volume (« Lobr », v. 419--420) ; la lecture est sûre ici par « Par ». Au-dessus du Q initial du paragraphe suivant, une petite marque en forme d'astérisque.}

Quand on demande le mouvement dela corde après qu'elle aura reçu
une impulsion quelconque, il faut absolument que la figure qui lui a été
imprimée d'abord soit telle, que non seulement toutes les ordonnées soient
presqu'infiniment petites mais que l'inclinaison de tous les élémens
de la courbe soient aussi infiniment petite. On pourroit encore ajouter
cette condition qu'on ait imprimé en même tems a chaque élement dela
corde un certain mouvement selon la direction de l'ordonnée et ce
mouvement initial doit aussi etre telle, qu'il n'en resulte aucun
saut dans la suite ; ou bien que la figure de la corde demeure
toujours conforme aux loix prescrites.

\emph{Question}

Ayant reduit la corde tendue à une figure quelconque, si aumoment
qu'on la relache, on imprime encore a chaque élement de la corde un
mouvement quelconque ; on demande pour chaque moment du tems suivant
tant la figure que le mouvement de la corde.

Soit $AB$ la corde fixe dans ses deux extremités $A$ et $B$, et tendue par une
force quelconque, a laquelle on ait imprimé aucommencement la figure
$ASB$, et d'abord cette courbe doit
être telle que 1.\textsuperscript{o} toutes ses appliquées
soient presqu'infiniment petites et 2.\textsuperscript{o}
que toutes ses tangentes ne s'éloignent
qu'infiniment peu de l'axe $AB$.
\note{« Question » est centré et souligné ; un court trait le sépare de ce qui suit. À droite de ces quatre lignes, une figure à la plume : l'axe horizontal $AB$ ; au-dessus et au-dessous, une grande courbe fermée qui passe par $A$ et $B$, coupée en $M$ (en haut) et $V$ (en bas) par une verticale qui rencontre l'axe en $X$ ; entre l'axe et la branche supérieure, une courbe plus plate, de $A$ à $B$, qui coupe la verticale en $S$, juste au-dessus de $X$.}

Ces deux conditions sont si naturellement liées avec la tension de la
corde a une telle figure ou ces deux conditions n'eussent pas lieu, dela
il est clair que la figure initiale peut être variée a l'infini, et qu'elle
dépend entièrement de nôtre volonté. Il est donc possible de donner
a la courbe une telle figure, que ne sauroit être exprimée par aucune
équation analytique, comme si on la tiroit par un mouvement libre
de la main, sans qu'aucune loi de continuité y ait lieu.
\note{La première phrase est écrite ainsi ; elle ne se construit pas, et rien ne semble y manquer sur la feuille.}

Il n'y certainement aucun doute, qu'on ne puisse imprimer a la corde
une telle figure, et ou pourtant les deux conditions aient lieu. Pour s'en
assurer on n'a qu'a tirer de $AOB$ une ligne courbe quelconque $AMB$,
en observant cette seule condition, qu'il n'y ait nulle part une tangente
perpendiculaire a l'axe : alors en diminuant toutes les appliquées $XM$
presqu'a l'infini selon un même rapport desorte que $XS=\alpha XM$,
prenant $\alpha$ pour une fraction extrèmement petite, non seulement toutes les
appliquées $XS$ deviendront infiniment petites, mais aussi les tangentes dans tous les points
$S$ seront infiniment peu inclinées a l'axe $AB$, tout comme les deux conditions prescrites l'exigent :
\note{« n'y certainement » est écrit sans verbe. « $AOB$ » : la lettre du milieu est un petit rond, qui ne figure pas sur le dessin. Le coefficient rendu $\alpha$ est un petit « u » ; les deux dernières lignes, plus serrées, vont jusqu'au bord du feuillet. Le texte continue au dos (vue 741, hors de ce lot).}

\end{document}
